\part{\trPicardLefschetzTheory}
\markboth{III　\trPicardLefschetzTheory}{III　\trPicardLefschetzTheory}
我们这套材料中的第三种成分，是\trLefschetzFibration{}的\trSymplectic{}几何。这个论题的内容异常广博，我们甚至远远谈不上把它讲得公允。此处使用的定义与 [131] 中的基本相同。事实上，我们的许多讨论都是对那篇论文的引申或评注（这意味着我们的处理颇为片面；想要获得更全面图景的读者，不妨把 [40]、[64]、[11]、[17] 作为可能的出发点）。

头两节回顾一些初等的方面。相对于 [131] 的主要补充是\trMatchingCycle{}的概念，它源出\trDonaldson{}（另见 [16]）。粗略地说，当我们有一个 \trLefschetzFibration，其一般\trFiber{}本身又是次级\trLefschetzFibration{}的\trTotalSpace{}时，\trMatchingCycle{}便出现了。我们把这样的构造称为\trBifibration；在最低的非平凡维数下，它是 [12]、[13] 中所考虑的辛\trBranchedCover{}的一个不那么精巧的版本。

在那之后，我们考虑\trFloer{} TQFT 的一个扩张，它基于\trLefschetzFibration{}的\trHolomorphic\trSection，并用它来推导一个基本的比较定理，连接 §5 的代数理论与 §16 的辛几何。这个定理在\trObject{}的层面上把\trTwistFunctor{} $T$ 的作用与一个\trSymplectic{}\trDehnTwist{} $\tau$ 的作用等同起来；它实际上是 [131] 中那条\trLongExactSequence{}的重新表述。我们不会重演那条序列的完整证明，但会较为详细地解释如何把这个结果变换成此处所需的形式。

给定一个以\trDisk{}为底空间的\trLefschetzFibration{}$\pi\colon E\to D$，可以取沿一根\trFiber{}$E_z$ 分歧的\trDoubleCover（在代数几何中，这会被称为“基变换”（base extension））；实际上，让 $z$ 变动就得到这样的\trCovering{}的一个\trFamily，而这正是\trBifibration{}的一个特别简单的例子。这个\trDoubleCover{}显然带有 $\ZZ/2$-作用，于是可以把 §6 与 §14 的\trEquivariant{}技巧施加于其上。我们将利用这个想法来定义一个 \trACategory{}$\scr{F}(\pi)$，称为这个\trLefschetzFibration{}的\trFukayaCategory；它介于\trFiber{}$E_z$ 的\trFukayaCategory{}与\trTotalSpace{}$E$ 的\trFukayaCategory{}之间，起中介的作用。无可否认，经由\trDoubleCover{}技巧的定义并不是最自然的一个，但它使人们能够重复利用大量现成的机器。

接着我们考察古典代数几何意义下的\trLefschetzPencil，并证明引言中所陈述的那些结果。最后，最后一节致力于一列例子，它们源自代数几何中一个稍有不同的部分，即 $(A_m)$ 型简单\trSingularPoint{}的\trMilnorFibre。这些例子主要具有历史性的意义，因为它们是最早提出导出\trFukayaCategory{}的显式代数模型的情形之一 [85]（事后来看，这些模型只在实维数 $\ge 4$ 时才是正确的）。到了那一步，讨论就变得不如本书其余部分那样自足，但有希望展示出在具体计算中所使用的那类几何与代数技巧。至于同一套一般性机器的更复杂的应用，我们请读者查阅文献，例如 [14]、[15]、[130]、[59]。
